%=====================================================================
\chapter{Securing the Shop}\label{ch:security}
%=====================================================================
\locator{4}{Part IV --- Money and Identity}

\begin{outcomes}
By the end of this chapter you will be able to:
\begin{tight}
  \item \textbf{State} the legal boundary of security testing, and
        \textbf{confine} your own work inside it.
  \item \textbf{Measure} a shop's security posture as a set of yes/no
        questions, and \textbf{repeat} the measurement after a change.
  \item \textbf{Distinguish} faults the framework already prevents from faults
        the deployment introduces, with evidence for each.
  \item \textbf{Apply} a hardening configuration, and \textbf{quantify} what it
        closed and what it did not.
  \item \textbf{Judge} which remaining findings need a decision rather than a
        setting, and \textbf{argue} for one.
\end{tight}
\end{outcomes}

\begin{prereq}
\chref{stack}: the container, Apache, and where configuration lives.
\chref{variations}: the upload directory and what is served from it.
\chref{accounts}: passwords, lockout, and that a defence can create a leak.
\textbf{Appendix~A's scope statement, which you must read before the lab.}
\end{prereq}

\begin{keyterms}
threat model $\bullet$ attack surface $\bullet$ security header $\bullet$
content security policy $\bullet$ clickjacking $\bullet$ MIME sniffing
$\bullet$ CORS $\bullet$ information disclosure $\bullet$ hardening $\bullet$
defence in depth $\bullet$ SQL injection $\bullet$ cross-site scripting
$\bullet$ CSRF token $\bullet$ PECA 2016
\end{keyterms}

\section{Where You May Point This Chapter}

\begin{alertbox}[Read this before the laboratory. It is not a formality.]
Every exercise in this chapter runs against \textbf{the container on your own
machine}, reached at \texttt{localhost}, which you installed, which contains no
real person's data, and which you can destroy and rebuild with one command.

\smallskip
\textbf{Running the same checks against any system you do not own and have not
been authorised in writing to test is a criminal offence in Pakistan} under the
\textbf{Prevention of Electronic Crimes Act 2016}, and under the equivalent law
of almost every other country. Sections 3 and 4 of that Act cover unauthorised
access to an information system and to data, with imprisonment and fine
attached to each.

\smallskip
There is no educational exemption. ``I was learning'', ``I did no harm'', ``the
site was insecure'' and ``I told them afterwards'' are not defences, and a
university assignment is not authorisation. The authorisation that counts is a
written scope agreement from the system's owner, naming what may be tested and
when.

\smallskip
\textbf{What this means in practice:} the only address that appears in this
chapter is \texttt{localhost:8090}. If you find yourself typing anything else,
stop.
\end{alertbox}

\section{Measure First, and Measure the Right Thing}

The temptation in a security chapter is to start with the famous attacks. It is
worth resisting, because the famous attacks are the ones the framework already
handles, and starting there teaches you to look in the wrong place.

\begin{measured}[The two attacks everyone expects, against your own search box]
From \texttt{code/m16\_security.py --injection}. Both probes are harmless: one
is a quotation mark, the other is a script tag that would open an alert box in
your own browser.

\rowcolors{2}{white}{lightbg}
\begin{longtable}{@{}L{5.2cm}L{2.2cm}L{2.4cm}L{2.2cm}@{}}
\toprule
\rowcolor{primary!15}
\textbf{Probe} & \textbf{Bytes} & \textbf{Verbatim?} & \textbf{Escaped?} \\
\midrule
\endhead
a normal search & 38{,}406 & yes & yes \\
one quotation mark & 26{,}440 & no & no \\
\texttt{ipod' OR '1'='1} & 26{,}508 & yes & yes \\
\texttt{<script>alert(1)</script>} & 26{,}648 & \textbf{no} & yes \\
\texttt{<img src=x onerror=\dots>} & 26{,}626 & \textbf{no} & yes \\
\bottomrule
\end{longtable}

\smallskip
\textbf{The SQL probe returns a \emph{shorter} page than the plain search}
---~26{,}508 bytes against 38{,}406. It found nothing, because the whole string
including the quotation marks was searched for literally. The query's meaning
did not change, which is what a parameterised or properly escaped query means.

\smallskip
\textbf{The script probes never appear as markup}, only in their escaped form.
Twig auto-escapes everything it prints unless explicitly told not to, so the
tag becomes text.

\smallskip
\textbf{So both classic attacks are already handled.} That is the useful
finding, and it redirects the rest of the chapter: whatever is wrong with this
shop is not in its application code.
\end{measured}

\section{What Is Actually Wrong, and It Is All Deployment}

\begin{measured}[Fifteen questions, asked of the shop as installed]
From \texttt{code/m16\_security.py}. Each is a yes/no question you could check
by hand with \texttt{curl}. ``Pass'' always means the safe answer.

\textbf{Part one: what the shop says about itself in every response.}

\rowcolors{2}{white}{lightbg}
\begin{longtable}{@{}L{7.4cm}L{1.9cm}L{1.9cm}L{2.4cm}@{}}
\toprule
\rowcolor{primary!15}
\textbf{Question} & \textbf{Before} & \textbf{After} & \textbf{Cost to fix} \\
\midrule
\endhead
Server header hides the version & fail & pass & 1 line \\
\texttt{X-Powered-By} absent & fail & pass & 1 line \\
\texttt{X-Content-Type-Options: nosniff} & fail & pass & 1 line \\
\texttt{X-Frame-Options} or frame-ancestors & fail & pass & 1 line \\
Content-Security-Policy set & fail & pass & 1 block \\
\texttt{Referrer-Policy} set & fail & pass & 1 line \\
\texttt{Permissions-Policy} set & fail & pass & 1 line \\
No wildcard CORS with credentials & fail & pass & 2 lines \\
\bottomrule
\end{longtable}

\smallskip
\textbf{Part two: what the shop will hand over if you ask for it.}

\rowcolors{2}{white}{lightbg}
\begin{longtable}{@{}L{7.4cm}L{1.9cm}L{1.9cm}L{2.4cm}@{}}
\toprule
\rowcolor{primary!15}
\textbf{Question} & \textbf{Before} & \textbf{After} & \textbf{Cost to fix} \\
\midrule
\endhead
\textbf{Error log not readable over HTTP} & \textbf{fail} & \textbf{pass} &
1 block \\
Session cookie is HttpOnly & pass & pass & --- \\
Session cookie has SameSite & pass & pass & --- \\
\texttt{system/storage/} not browsable & pass & pass & --- \\
\texttt{install/} removed & pass & pass & --- \\
\texttt{config.php} not served as text & pass & pass & --- \\
Admin not at the default path & fail & \textbf{fail} & a decision \\
\midrule
\textbf{Both parts together} & \textbf{5 of 15} & \textbf{14 of 15} & \\
\bottomrule
\end{longtable}

\smallskip
\textbf{One configuration file closed nine findings.} Not one of them required
a change to the shop's code, and not one of them took more than a few lines.

\smallskip
\textbf{The serious one is the error log.} Before hardening, anybody at all who
asked the shop for \texttt{system/storage/logs/error.log} was given it ---
verified by writing a marker string into the file and then fetching it back
over HTTP. A PHP error log contains file paths, SQL fragments, and often the
contents of whatever request failed.
\end{measured}

\begin{alertbox}[Why the shipped protection does not work, in two parts]
The platform knows about this. It ships a file called
\texttt{.htaccess\textbf{.txt}} containing, among other things, exactly the rule
that would have blocked the log. Every guide to securing this software says:
rename it.

\smallskip
\textbf{First, nobody renames it.} The installer does not, the Docker image does
not, and a fresh install therefore has no \texttt{.htaccess} at all.

\smallskip
\textbf{Second, and worse: renaming it would not help here.} Debian's
\texttt{apache2.conf} sets \texttt{AllowOverride None} for the document root,
so Apache never reads \texttt{.htaccess} files. You could rename it, confirm
the file is in place, and have changed nothing --- with no error, no warning,
and no way to tell.

\smallskip
\textbf{This is why the fix in \texttt{code/shop/security.conf} is a server
configuration rather than a \texttt{.htaccess}.} It is also why the advice you
will find online for this platform is worse than useless: it is confidently
wrong, it has been repeated for fifteen years, and following it produces a
shop that is exactly as exposed as before while its owner believes otherwise.

\smallskip
\textbf{Verify that a control works. Do not verify that you installed it.}
\end{alertbox}

\section{Three Mistakes Made While Writing This Chapter}

The hardening file did not work the first time, or the second. The failures are
kept because each one is a way a real control silently does nothing.

\begin{examplebox}[What went wrong, and how the measurement found it]
\begin{tightnum}
  \item \textbf{\texttt{Header always unset} removed nothing.} The CORS headers
        were still there after hardening. \texttt{mod\_headers} keeps two
        tables: \texttt{always} applies to every response including errors, and
        the default applies only to successful ones. A header PHP sets on a 200
        lives in the second table. \textbf{Both forms are needed}, and using
        only the one whose name sounds more thorough removes nothing.
  \item \textbf{\texttt{ServerTokens} was ignored after a reload.} It is read
        only when Apache \emph{starts}. A graceful reload applies every other
        directive in the file and silently skips this one.
  \item \textbf{The file was overridden by Debian's own.} Apache reads
        \texttt{conf-enabled} in filename order, and the last setting of a
        server-level directive wins. A file named
        \texttt{security-ecad.conf} sorts \emph{before}
        \texttt{security.conf} --- because \texttt{-} precedes \texttt{.} ---
        so Debian's \texttt{ServerTokens OS} was applied after it. Renaming
        the file to \texttt{zz-security-ecad.conf} fixed it.
\end{tightnum}

\smallskip
\textbf{All three produced a valid configuration that Apache accepted without
complaint, and in all three the header was simply absent.} A security control
that fails open and fails quietly is the only kind worth being afraid of, and
the only defence is to measure the result rather than the action. Each of these
took one run of the harness to find and would have survived any amount of
re-reading the file.
\end{examplebox}

\dsfig{%
\begin{tikzpicture}[
  lyr/.style={rounded corners=3pt, draw=#1, line width=1.1pt, fill=#1!8,
              text=#1!60!black, font=\bfseries\small, align=center,
              text width=3.0cm, minimum height=1.05cm},
  note/.style={font=\scriptsize, text=neutral, align=left, text width=7.6cm,
               anchor=west},
  tick/.style={font=\scriptsize\bfseries, text=greenhl},
  cross/.style={font=\scriptsize\bfseries, text=accent}
]
  \node[lyr=greenhl]   (app) at (0,3.3) {Application\\{\scriptsize\mdseries
       the shop's own code}};
  \node[lyr=accent]    (svr) at (0,1.65) {Server\\{\scriptsize\mdseries
       Apache configuration}};
  \node[lyr=purplehl]  (dec) at (0,0)    {Decisions\\{\scriptsize\mdseries
       how it is deployed}};

  \node[note] at (2.0,3.3) {\textcolor{greenhl}{\faIcon{check}}~Escapes SQL.
       \textcolor{greenhl}{\faIcon{check}}~Auto-escapes output.
       \textcolor{greenhl}{\faIcon{check}}~CSRF tokens on forms.
       \textcolor{greenhl}{\faIcon{check}}~bcrypt passwords.
       \emph{Nothing to fix here.}};
  \node[note] at (2.0,1.65) {\textcolor{accent}{\faIcon{times}}~Nine findings,
       every one a line or two of configuration --- headers, version
       disclosure, the readable error log. \emph{One file fixed all of them.}};
  \node[note] at (2.0,0)   {\textcolor{accent}{\faIcon{times}}~Admin on the
       default path; storage inside the web root; no HTTPS.
       \emph{These need somebody to decide, not to type.}};
\end{tikzpicture}}%
{Where a stock shop's security problems actually live. The layer students
expect to fix is the one that was already correct.}{securitylayers}

\section{What a Header Cannot Fix}

One finding survived the configuration file, and it is the honest one.

\begin{conceptbox}[The three that need a decision]
\begin{tightnum}
  \item \textbf{The admin panel is at \texttt{/admin/}.} Renaming the directory
        and the matching path in \texttt{admin/config.php} is not security ---
        anyone who looks will find it --- but it removes the shop from every
        automated scan that tries one path and moves on. It costs a
        configuration change and a note for whoever maintains the shop next,
        which is why it is a decision rather than a setting.
  \item \textbf{\texttt{system/storage/} is inside the web root.} Blocking it
        in the server config, as this chapter did, is the fix you can apply
        today. \textbf{Moving it outside the web root is the fix that is still
        correct after somebody edits the server config in two years.} Prefer
        the control that does not depend on remembering.
  \item \textbf{There is no HTTPS.} On \texttt{localhost} it does not matter.
        On a real shop everything in this book --- the session cookie, the
        password, the order --- travels in clear text without it, and
        \texttt{Strict-Transport-Security} is deliberately left out of the
        hardening file because switching it on without a certificate makes the
        shop unreachable in your own browser for months.
        \chref{deploy} turns it on where there is something to turn it on for.
\end{tightnum}
\end{conceptbox}

\begin{alertbox}[The defaults this course has been using, and must not keep]
The shop was installed with \texttt{admin}/\texttt{admin123}, a database user
\texttt{shop}/\texttt{shoppw}, and a root password of \texttt{rootpw}. Those are
in the handout, in the repository, and in this paragraph.

\smallskip
They are acceptable for exactly one thing: a container on your own machine,
with no data in it, not reachable from your network. They are documented as
teaching defaults rather than hidden, because a weak password nobody mentions
is worse than a weak password everybody knows about.

\smallskip
\textbf{Any shop that leaves this machine changes all three before it does},
along with the minimum password length of six characters noted in
\chref{accounts}. \chref{deploy} makes that a checklist item with nothing
optional about it.
\end{alertbox}

\begin{pitfall}
\begin{tight}
  \item \textbf{Testing anything but your own container.} See the first page of
        this chapter, and Appendix~A.
  \item \textbf{Starting with SQL injection and cross-site scripting.} Both
        were already handled; the time was better spent elsewhere.
  \item \textbf{Renaming \texttt{.htaccess.txt} and believing you are done.}
        \texttt{AllowOverride None} means the file is never read.
  \item \textbf{Verifying that you applied a control rather than that it
        works.} Three separate controls in this chapter were applied correctly
        and did nothing.
  \item \textbf{Using \texttt{Header always unset} alone} to remove a header
        from a successful response.
  \item \textbf{Reloading instead of restarting} after changing a directive
        that is only read at startup.
  \item \textbf{Setting \texttt{Strict-Transport-Security} without a
        certificate}, which locks you out of your own shop.
  \item \textbf{Writing a Content-Security-Policy that breaks the theme}, then
        removing it entirely instead of narrowing it.
  \item \textbf{Treating a passing checklist as a secure shop.} Fourteen of
        fifteen is a measurement, not a guarantee.
\end{tight}
\end{pitfall}

\begin{lab}[Break into your own shop, then shut the door]
\begin{lstlisting}[language=bash]
# SCOPE: localhost only. Your container, your machine, no real data.
# Re-read the first page of this chapter before you begin.

# --- 1. Ask what it tells you, before changing anything -------------
cd code
curl -sI http://localhost:8090/ | head -20
# List every header that names a version or a product.

# --- 2. Read your own error log over HTTP ---------------------------
docker compose exec shop sh -c \
  'echo "SECRET-TEST-STRING" >> system/storage/logs/error.log'
curl -s http://localhost:8090/system/storage/logs/error.log
# If you see your string, so would anyone else.

# --- 3. Run the battery ---------------------------------------------
python m16_security.py --check

# --- 4. Try the two attacks you came here for -----------------------
python m16_security.py --injection
# Predict the result first. Most people predict at least one works.

# --- 5. Find out why the shipped protection does not ----------------
docker compose exec shop sh -c 'ls -la .htaccess*'
docker compose exec shop sh -c \
  'grep -n AllowOverride /etc/apache2/apache2.conf'
# Two separate reasons. Write both down.

# --- 6. Harden, and measure the difference --------------------------
python m16_security.py
# Read shop/security.conf first and predict the score.

# --- 7. Check the shop still works ----------------------------------
python waitfor.py
python m10_checkout.py --repeat 1
python m10_checkout.py --clean
# A hardening change that breaks checkout is not a hardening change.

# --- 8. Close the last finding --------------------------------------
# Rename the admin directory AND the path in admin/config.php. Then
# re-run the battery. Say why this is weaker than the other fixes.

# --- 9. Write the five lines you would hand to a client -------------
# What you changed, what is still open, what each one would cost,
# and which one you would do first.
\end{lstlisting}

\smallskip
\textbf{What to notice.} Step 4 against step 3. The attacks you have heard of
were already stopped; the shop was giving away its error log. Security attention
goes where the stories are, and the stories are about the wrong layer.
\end{lab}

\begin{checkpoint}
\begin{tightnum}
  \item A colleague renames \texttt{.htaccess.txt} to \texttt{.htaccess} on
        this stack and reports the shop secured. Give the two independent
        reasons that is wrong, say what you would run to demonstrate it in one
        command, and state the general lesson about verifying controls.
  \item The battery went from 5 of 15 to 14 of 15 with one configuration file.
        Explain why none of the nine fixes touched the shop's code, what that
        says about where to look when you inherit a shop, and which single
        remaining finding you would address first and why.
  \item Three controls in this chapter were written correctly, accepted by
        Apache without complaint, and did nothing. Name the mechanism behind
        any two, and describe a working habit that would catch all three
        without knowing any of them in advance.
\end{tightnum}
\end{checkpoint}

\begin{chaptersummary}
\begin{tight}
  \item Every exercise is confined to your own container on \texttt{localhost}.
        Testing anything else is a criminal offence under PECA 2016, with no
        educational exemption.
  \item Measured: SQL injection and cross-site scripting against the search box
        were both already handled --- the injection string was searched for
        literally and returned fewer results, and the script tag appeared only
        escaped. The famous attacks were not the problem.
  \item Measured: the shop as installed passed 5 of 15 checks. After one
        configuration file it passed 14. Nothing in the shop's code changed.
  \item The serious finding was the PHP error log, readable over HTTP ---
        confirmed by writing a marker into it and fetching it back.
  \item The shipped protection fails twice over: \texttt{.htaccess.txt} is
        never renamed, and \texttt{AllowOverride None} means it would not be
        read if it were. The advice repeated everywhere online for this
        platform does nothing on a modern deployment.
  \item Three controls were applied correctly and silently did nothing:
        \texttt{Header always unset} on a successful response,
        \texttt{ServerTokens} after a reload rather than a restart, and a file
        that sorted before Debian's own. Verify the result, never the action.
  \item What a header cannot fix needs a decision: move the admin path, move
        storage out of the web root, and serve over HTTPS. Prefer the control
        that does not depend on somebody remembering.
  \item The default credentials used throughout this course are documented
        rather than hidden, and are acceptable for a local container and
        nothing else.
\end{tight}
\end{chaptersummary}

\begin{reviewq}
  \item \textbf{(MCQ)} The stock shop passed how many of the fifteen checks?
        (a)~0; (b)~5; (c)~11; (d)~14.
  \item \textbf{(MCQ)} Renaming \texttt{.htaccess.txt} to \texttt{.htaccess}
        on this stack achieves: (a)~full protection; (b)~partial protection;
        (c)~nothing, because Apache does not read it; (d)~nothing, because the
        file is empty.
  \item \textbf{(Short)} Explain what \texttt{X-Content-Type-Options: nosniff}
        prevents, and connect it to \chref{variations}'s upload findings.
  \item \textbf{(Short)} Why is \texttt{Strict-Transport-Security} deliberately
        absent from this chapter's hardening file?
  \item \textbf{(Short)} Give one reason moving \texttt{system/storage/} out of
        the web root is better than blocking it in the server configuration,
        even though both work today.
  \item \textbf{(Applied)} You inherit a shop you did not build. Write the
        order in which you would investigate it, with the specific command or
        observation at each step, and say at which point you would stop and
        report rather than continue --- and why that point exists.
  \item \textbf{(Applied)} Design a Content-Security-Policy for a shop whose
        theme uses inline scripts, a payment provider's hosted field, and an
        analytics service. State each directive, what it allows, what it
        blocks, how you would roll it out without breaking the shop, and how
        you would know it was working.
\end{reviewq}
